\documentclass[runningheads]{llncs}

\usepackage[T1]{fontenc}
\usepackage{newtxtext}
\usepackage[varvw]{newtxmath}
\usepackage{graphicx}
\usepackage{booktabs}
\usepackage{orcidlink}

\newcommand{\todo}[1]{\textbf{[TODO: #1]}}
\newcommand{\TBD}{\textbf{TBD}}
\newcommand{\mstd}[2]{#1{\scriptsize$\pm$#2}}

\begin{document}

\title{Beyond In-Domain Discrimination: A Multi-Source Study of Identity-Learning Objectives for Cross-Domain Person Re-Identification}
\titlerunning{Identity-Learning Objectives for Person Re-ID}

\author{Ahmed Labib\inst{1}\textsuperscript{*}\textsuperscript{\ensuremath{\dagger}}\,\orcidlink{0009-0002-6835-4422} \and
Mohamed ElBassat\inst{1}\textsuperscript{*}\textsuperscript{\ensuremath{\dagger}}\,\orcidlink{0009-0006-3917-1319} \and
Mohamed Ehab\inst{1}\textsuperscript{\ensuremath{\dagger}}\,\orcidlink{0009-0001-0507-8099} \and
Rokayya Aly\inst{1}\,\orcidlink{0009-0001-6002-283X} \and
Ahmed Temraz\inst{1}\,\orcidlink{0009-0001-1632-2704} \and
Malak Alaaeldin\inst{1}\,\orcidlink{0009-0003-2737-4045} \and
Omar Nabil\inst{1}\,\orcidlink{0009-0003-7528-0131} \and
Yara Hossam\inst{1}\,\orcidlink{0009-0001-7838-1072} \and
Ola Omara\inst{1}\,\orcidlink{0009-0007-3602-875X} \and
Jana Wael\inst{1}\,\orcidlink{0009-0002-5751-4144} \and
Mahmoud Gamal\inst{1}}

\authorrunning{A. Labib et al.}

\institute{Faculty of Computers and Data Science, Alexandria University, Alexandria, Egypt\\
\textsuperscript{*} Correspondence: \email{cds.AhmedMahmoud24052@alexu.edu.eg} (A. Labib); \email{cds.MohamedMohamed24510@alexu.edu.eg} (M. ElBassat).\\
\textsuperscript{\ensuremath{\dagger}} These authors contributed equally.}

\maketitle

\begin{abstract}
Person re-identification (Re-ID) models are often developed and evaluated within a single benchmark family, so their performance on familiar datasets may not predict transfer to a new camera network. Multi-source training broadens capture conditions but also exposes models to source-dependent annotation quality and substantial within-identity variation. We ask whether the identity-learning objective affects generalization to unseen domains. Using the FastReID Strong Baseline with a ResNet-50-IBN backbone, we first compare eight candidate configurations: four identity-learning heads (linear softmax, ArcFace, Sub-center ArcFace, and SphereFace2), each trained with and without triplet loss. Each configuration is trained with three random seeds and evaluated on four unseen public benchmarks and a pilot surveillance-style campus set from Alexandria University. We then compare the candidate selected by internal validation against the FastReID Circle-softmax-plus-triplet default across the Strong Baseline, OSNet, and ConvNeXt-S. Because the identity classifier is discarded at inference, changing the objective adds no classification-head inference cost within a backbone; the cross-backbone comparison assesses architecture-dependent tradeoffs. \todo{Report the selected candidate, its gains over the FastReID default, seed variation, and Alexandria performance.}
\end{abstract}

\keywords{Person re-identification \and Identity-learning objectives \and Domain generalization \and Multi-source training \and Angular-margin losses}

%==============================================================================


\section{Introduction}
\label{sec:intro}

Person re-identification (Re-ID) aims to retrieve images of the same individual across non-overlapping camera views. Despite advances in deep representation learning, the task remains challenging due to variations in viewpoint, illumination, occlusion, resolution, and clothing. Convolutional architectures have established effective and reproducible Re-ID baselines through learned identity representations and metric-learning objectives \cite{luo2019bag,sun2018beyond,song2020strong}. Compact designs such as OSNet further demonstrate their relevance to resource-constrained applications \cite{zhou2022learning}. Although Transformer-based and vision-language approaches have advanced benchmark performance \cite{he2021transreid,li2023clip,chen2023beyond}, CNN-based pipelines remain attractive when computational resources and model complexity are practical constraints.

An important challenge in deploying Re-ID systems is generalization to camera environments not represented during training. Differences in acquisition conditions and appearance distributions can cause substantial domain shifts, making strong in-domain performance an unreliable indicator of performance on unseen camera networks \cite{jin2020snr,zhou2021mix,zhao2021m3l,lee2025domain}. Multi-source training can improve exposure to diverse conditions, but combining heterogeneous datasets also introduces differences in annotation quality, sampling distributions, and within-identity appearance variation, particularly in datasets containing clothing changes, occlusion, and video-derived observations \cite{shu2022last,xu2023deepchange}.

Beyond the training data and feature extractor, the identity-learning objective determines how the embedding space is structured. Classification objectives encourage identity discrimination, while metric-learning objectives such as triplet loss constrain relative distances between samples \cite{hermans2017triplet}. Alternative classification heads introduce different assumptions about class separation and intra-class variation, including angular margins, multiple class centers, and independent binary classification \cite{deng2019arcface,deng2020subcenter,wen2022sphereface2}. However, the relative behavior of these objectives under heterogeneous multi-source training and evaluation across unseen domains remains insufficiently characterized. This motivates a controlled investigation in which objective selection is examined independently of the feature-extraction architecture.

Accordingly, we first compare eight objective configurations using the FastReID Strong Baseline (SBS): four identity-learning heads, each evaluated with and without triplet loss, under a common training protocol and across three random seeds. The FastReID Circle-softmax-plus-triplet default is excluded from this loss sweep. After selecting the strongest candidate using composite internal validation, we compare it with that default across the SBS, OSNet, and ConvNeXt-S. Evaluation uses held-out Re-ID benchmarks and an independently collected pilot dataset from Alexandria University. Since the identity-classification head is discarded during retrieval inference, changing the training objective adds no classification-head computations at inference.

\paragraph{Contributions.}
This work contributes a controlled, three-seed comparison of eight identity-learning objective configurations, covering four heads with and without triplet supervision under heterogeneous multi-source training. It also compares the candidate selected from this loss sweep with the FastReID Circle-softmax-plus-triplet default across the FastReID Strong Baseline, OSNet, and ConvNeXt-S. In addition, the study establishes a harmonized nine-source Re-ID training protocol with common metadata representation, domain-balanced sampling, composite internal validation, and held-out evaluation, and reports an independent pilot evaluation on surveillance-style Re-ID data collected at Alexandria University.


%==============================================================================

\section{Related Work}
\label{sec:related}

\paragraph{Domain-generalizable person re-identification.}
Domain generalization in person Re-ID aims to learn identity representations from labeled source domains that remain effective on unseen camera networks without target-domain data. Prior approaches address this challenge through style normalization, feature-statistics mixing, and multi-source training strategies. Style Normalization and Restitution (SNR) suppresses style variations through instance normalization while restoring identity-relevant information~\cite{jin2020snr}; MixStyle mixes feature statistics across training samples to synthesize style variation~\cite{zhou2021mix}; and Memory-based Multi-Source Meta-Learning (M$^3$L) simulates unseen-domain shifts with meta-learning, a memory-based identification loss, and meta-batch normalization~\cite{zhao2021m3l}. Lee et al.'s 2025 survey reviews these and related domain-generalization strategies~\cite{lee2025domain}. Whereas these methods modify representations or training procedures, our work isolates the identity-learning objective under a shared multi-source protocol.

\paragraph{Identity-learning objectives.}
Person Re-ID commonly employs identity classification and metric-learning objectives, either independently or jointly. Triplet loss encourages relative distance separation between samples of the same and different identities \cite{hermans2017triplet}. Angular-margin approaches, including SphereFace and ArcFace, introduce geometric constraints to improve inter-class discrimination \cite{liu2017sphereface,deng2019arcface}. Sub-center ArcFace extends this formulation by representing each identity with multiple centers, providing greater flexibility for intra-class variation and noisy labels \cite{deng2020subcenter}. SphereFace2 adopts an independent binary-classification formulation motivated by open-set recognition \cite{wen2022sphereface2}, whereas Circle-based objectives adapt the weighting of similarity relationships during optimization \cite{sun2020circle,song2020strong}. Recent work has also examined alternative loss combinations for cross-domain Re-ID. AATransID, for instance, evaluates combinations of ArcFace, Circle, CosFace, and additive-margin softmax with triplet supervision, reporting its strongest results with ArcFace and triplet loss \cite{yildiz2026aatransid}. Unlike this Transformer-based approach, our study compares a broader set of objective configurations within a common CNN-based training framework and evaluates their behavior across multiple held-out domains.

\paragraph{Convolutional architectures for person Re-ID.}
Convolutional architectures remain an established foundation for practical person Re-ID. Strong baselines combine classification supervision, metric learning, and batch-normalization necks to produce discriminative retrieval embeddings \cite{luo2019bag,sun2018beyond,song2020strong}. FastReID provides a unified framework for implementing and evaluating these components \cite{he2023fastreid}, while OSNet introduces omni-scale feature learning to capture discriminative information at different spatial scales with an efficient architecture \cite{zhou2022learning}. Transformer-based alternatives, including TransReID and CLIP-ReID, have further expanded the architectural landscape through attention-based and vision-language representations \cite{he2021transreid,li2023clip}. Rather than introducing a new feature extractor, our work uses the FastReID Strong Baseline for the primary objective comparison and subsequently evaluates the selected formulation on OSNet and ConvNeXt-S to examine its transfer across different CNN architectures.

%==============================================================================


\section{Datasets}
\label{sec:data}

Three data partitions are used in this study: nine heterogeneous sources assembled for training and validation, and four public benchmarks reserved exclusively for unseen-domain evaluation. The overall experimental pipeline is illustrated in Fig.~\ref{fig:framework}.

\begin{figure}[t]
\centering
\fbox{\parbox{0.92\textwidth}{\centering\vspace{1.8em}
\textbf{Figure 1 (placeholder).}\\
Multi-source training $\rightarrow$ eight candidate loss configurations on the Strong Baseline $\rightarrow$ validation-based candidate selection $\rightarrow$ comparison with the FastReID default across the Strong Baseline, OSNet, and ConvNeXt-S $\rightarrow$ unseen public benchmarks.
\vspace{1.8em}}}
\caption{Overview of the experimental framework.}
\label{fig:framework}
\end{figure}

\subsection{Multi-Source Training Data}
\label{sec:train-sources}

Training data comprise nine Re-ID datasets selected to represent diverse acquisition conditions, including occlusion, clothing changes, and multi-camera video observations. Their respective image counts are summarized in Table~\ref{tab:train}. Together, the curated sources contain 782,064 images and 10,036 identities. The four public evaluation benchmarks are excluded from training and model selection.

\begin{table}[t]
\centering
\caption{Composition of the multi-source training dataset.}
\label{tab:train}
\begin{tabular}{l r}
\toprule
\textbf{Dataset} & \textbf{Images} \\
\midrule
CUHK03-NP~\cite{li2014deepreid,zhong2017reranking} & 5,892 \\
LaST~\cite{shu2022last} & 71,248 \\
Occluded-PoseTrack-ReID~\cite{somers2024kpr} & 14,266 \\
DeepChange~\cite{xu2023deepchange} & 60,916 \\
LPW~\cite{su2016deep} & 338,952 \\
ORBench RGB~\cite{zuo2025reid5o} & 21,564 \\
PRID2011~\cite{hirzer2011person} & 31,690 \\
RPIfield~\cite{zheng2018rpifield} & 179,996 \\
IUST PersonReID~\cite{moghaddam2025iust} & 57,540 \\
\midrule
\textbf{Total} & \textbf{782,064} \\
\bottomrule
\end{tabular}
\end{table}

\subsection{Data Harmonization and Partitioning}
\label{sec:train-curation}

The source datasets differ in directory structure, identity conventions, camera metadata, and annotation formats. Each source was audited and converted into a common metadata schema with consistent sample identifiers, identity keys, and camera fields.
Identity labels were remapped into a global index space containing 10,036 identities, with identity spaces kept separate across source datasets. Each source was partitioned into training and validation identities, using an approximately 80/20 split where applicable and respecting source-specific release protocols. Validation identities were excluded from optimization and reserved for checkpoint selection (Section~\ref{sec:training}). To reduce domination by larger sources, particularly video-derived datasets, domain-balanced sampling was used during training, with source weights proportional to the square root of their identity counts. (needs more explanation)


\subsection{Unseen Public Benchmarks}
\label{sec:public-test}

Cross-domain retrieval is evaluated on four public benchmarks excluded from training and model selection: Market-1501~\cite{zheng2015scalable}, MSMT17~\cite{wei2018person}, DukeMTMC-reID~\cite{zheng2017unlabeled}, and Occluded-Duke~\cite{miao2019occluded}. These benchmarks cover different evaluation conditions, from conventional still-image retrieval to large-scale multi-camera and occlusion-heavy scenarios. Their query and gallery sizes are summarized in Table~\ref{tab:public}.

\begin{table}[t]
\centering
\caption{Query and gallery sizes of the unseen public evaluation benchmarks.}
\label{tab:public}
\begin{tabular}{l r r}
\toprule
\textbf{Dataset} & \textbf{Query} & \textbf{Gallery} \\
\midrule
Market-1501~\cite{zheng2015scalable} & 3,368 & 19,732 \\
MSMT17~\cite{wei2018person} & 11,659 & 82,161 \\
DukeMTMC-reID~\cite{zheng2017unlabeled} & 2,228 & 17,661 \\
Occluded-Duke~\cite{miao2019occluded} & 2,210 & 17,661 \\
\bottomrule
\end{tabular}
\end{table}

Occluded-Duke is derived from DukeMTMC; this relationship is acknowledged when interpreting results.

%==============================================================================

\section{Methodology}
\label{sec:method}

\subsection{Model Architecture}
\label{sec:architecture}

This work compares three existing convolutional architectures rather than introducing a new model. The primary eight-objective loss sweep uses the FastReID Strong Baseline (SBS)~\cite{he2023fastreid,song2020strong}. FastReID serves as the implementation framework, while the SBS combines a ResNet-50 backbone with instance-batch normalization (IBN)~\cite{pan2018two}, generalized-mean (GeM) pooling, and a batch-normalization neck to produce a 2048-dimensional embedding. The cross-backbone comparison additionally includes OSNet~\cite{zhou2022learning} and ConvNeXt-S~\cite{liu2022convnet}. It compares the FastReID Circle-softmax-plus-triplet default against the candidate selected from the eight-configuration loss sweep. All architectures use the same multi-source training data, input resolution, optimization schedule, and evaluation protocols. Training is implemented in PyTorch using the FastReID codebase, extended with our multi-source data loader and domain-balanced sampler (Section~\ref{sec:train-curation}).

\subsection{Training Setup}
\label{sec:training}
The primary SBS loss sweep evaluates eight candidate configurations that differ in identity-learning objective and triplet supervision (Section~\ref{sec:objectives}). Each architecture is initialized using its corresponding ImageNet-pretrained weights. We use Adam with a base learning rate of $3.5\times10^{-4}$ and weight decay of $5\times10^{-4}$. Cross-entropy-based heads use label smoothing with $\epsilon=0.1$, while configurations incorporating triplet loss use hard mining with a zero margin.

Training proceeds for 50 \emph{virtual epochs}, each comprising 2,000 parameter updates sampled from the multi-source loader. This definition accommodates the combined training pool, which is substantially larger than a single pass over all images. The backbone remains frozen for the first 1,000 updates and is subsequently optimized jointly with the embedding head. The learning rate follows cosine decay after a 20-epoch delay, with linear warmup over the first 1,000 iterations and a minimum learning rate of $7\times10^{-7}$. Minibatches contain 128 images, with eight instances per identity, using domain-balanced sampling as described in Section~\ref{sec:train-curation}. Mixed-precision forward passes are used during training, whereas validation feature extraction is performed in full precision. Primary experiments are conducted using two GPU accelerators with data-parallel model replication. Each SBS objective is trained using three random seeds ($0$, $1$, and $2$), with all other hyperparameters held fixed. The separate cross-backbone comparison applies the same optimization and evaluation protocol to the FastReID Circle-softmax-plus-triplet default and the candidate selected from the primary loss sweep.

For checkpoint selection, we construct a composite validation protocol using source-specific identities, retaining official partitions where available and predefined cross-view or cross-context splits otherwise. The four public evaluation benchmarks are excluded from training, validation, and checkpoint selection. Every five virtual epochs, including the final epoch, we evaluate each validation protocol using L2-normalized cosine similarity. Gallery images sharing both identity and camera with a query are excluded. The selected checkpoint maximizes the unweighted mean validation score across protocols:

\begin{equation}
s_i =
\frac{\mathrm{Rank\text{-}1}_i+\mathrm{mAP}_i}{2}.
\end{equation}

The final checkpoint is selected according to the mean of $s_i$ across validation protocols.

\subsection{Identity-Learning Objectives in a Unified View}
\label{sec:objectives}

We compare different identity-learning heads using a common notation that distinguishes their target and non-target scores and the class structures they assume. The implementations follow the cited formulations; the objective configurations are summarized in Table~\ref{tab:unified}.

\paragraph{Notation and cross-entropy heads.}
Let $\mathbf{f}\in\mathbb{R}^{d}$ denote the embedding of an image with identity label $y$, and let $C=10{,}036$ denote the number of training identities. Linear softmax uses bias-free logits $z_j=\mathbf{w}_j^{\top}\mathbf{f}$, following the BNNeck design~\cite{luo2019bag}. For angular heads, features and class centers are $\ell_2$-normalized, such that

\begin{equation}
\cos\theta_j =
\frac{\mathbf{w}_j^{\top}\mathbf{f}}
{\|\mathbf{w}_j\|\|\mathbf{f}\|},
\end{equation}

where $s$ denotes the logit scale. Cross-entropy-based heads use label smoothing with $\epsilon=0.1$. The smoothed target distribution is defined as

\begin{equation}
q_j =
\begin{cases}
1-\epsilon, & j=y,\\
\frac{\epsilon}{C-1}, & j\neq y.
\end{cases}
\end{equation}

Given the head-specific logits $z_j$ and softmax probabilities $p_j$, the classification objective is

\begin{equation}
\mathcal{L}_{\mathrm{CE}}
=
-\sum_{j=1}^{C}q_j\log p_j,
\qquad
p_j=\frac{\exp(z_j)}{\sum_{k=1}^{C}\exp(z_k)}.
\end{equation}

SphereFace2 instead employs an independent binary-classification objective, described below.

\paragraph{Classification heads.}
Table~\ref{tab:unified} summarizes the compared heads. For Sub-center ArcFace, each identity is represented by $K$ centers. The target identity uses its best-matching center, while each non-target identity is represented by its maximum-similarity center. The Circle-softmax formulation corresponds to the class-level variant implemented in the Strong Baseline rather than the pair-based Circle Loss formulation.

\begin{table}[t]
\caption{Unified view of the compared identity-learning heads. $\theta_j$ denotes the angle between a feature and a class center, while $z_y$ and $z_j$ denote target and non-target logits for cross-entropy heads.}
\label{tab:unified}
\centering
\scriptsize
\setlength{\tabcolsep}{2pt}
\begin{tabular}{@{}p{1.8cm}p{2.8cm}p{2.7cm}p{2.8cm}p{1.5cm}@{}}
\toprule
Head & Target score & Non-target score & Class assumption & Settings \\
\midrule
Linear softmax~\cite{luo2019bag}
& $\mathbf{w}_y^{\top}\mathbf{f}$
& $\mathbf{w}_j^{\top}\mathbf{f}$
& One unnormalized vector per identity
& None \\

ArcFace~\cite{deng2019arcface}
& $s\cos(\theta_y+m)$
& $s\cos\theta_j$
& One center; additive angular margin
& $s=64$, $m=0.35$ \\

Sub-center ArcFace~\cite{deng2020subcenter}
& $s\cos(\theta_{y,k^*}+m)$
& $s\max_k\cos\theta_{j,k}$
& $K$ centers per identity
& $K=3$ \\

SphereFace2~\cite{wen2022sphereface2}
& $r(g(\cos\theta_y)-m)+b$
& $r(g(\cos\theta_j)+m)+b$
& Independent per-class binary decisions
& $r=40$, $m=0.4$ \\

Circle-softmax~\cite{sun2020circle,song2020strong}
& $s\alpha_p(\cos\theta_y-\Delta_p)$
& $s\alpha_n^j(\cos\theta_j-\Delta_n)$
& One center; adaptive per-logit weights
& $s=64$, $m=0.35$ \\
\bottomrule
\end{tabular}
\end{table}

For Circle-softmax, the adaptive weights are defined as

\begin{equation}
\alpha_p=[1+m-\cos\theta_y]_+,
\qquad
\alpha_n^j=[\cos\theta_j+m]_+,
\end{equation}

with $\Delta_p=1-m$ and $\Delta_n=m$. These weights are treated as constants during backpropagation. This is the class-level head used by FastReID, rather than the pair-based objective introduced in the original Circle Loss formulation~\cite{sun2020circle,song2020strong}.

SphereFace2 replaces the softmax over identities with $C$ independent binary classifiers~\cite{wen2022sphereface2}. We use its SphereFace2-C formulation with an additive margin and similarity adjustment:

\begin{equation}
g(z)=2\left(\frac{z+1}{2}\right)^t-1.
\end{equation}

The resulting objective is

\begin{equation}
\begin{split}
\mathcal{L}_{\mathrm{SF2}}={}&
\frac{\omega\lambda}{r}
\log\left(1+\exp\left[-\left(r(g(\cos\theta_y)-m)+b\right)\right]\right)
\\
&+
\frac{\omega(1-\lambda)}{r}
\sum_{j\ne y}
\log\left(1+\exp\left[r(g(\cos\theta_j)+m)+b\right]\right).
\end{split}
\end{equation}

Here, $r$ denotes the hypersphere radius, $m$ the margin, $\lambda$ balances the positive term against the negative terms, $b$ is a learnable bias, and $\omega$ is a constant loss weight. We adopt the reference implementation settings $r=40$, $m=0.4$, $\lambda=0.9$, $t=3$, and $\omega=50$. The bias $b$ is initialized to balance the positive and negative terms at the beginning of training. As in the reference implementation, the similarity adjustment and margin are treated as constant offsets during backpropagation, with gradients computed with respect to the cosine similarities.

\paragraph{Triplet loss and evaluated configurations.}
When enabled, hard-mined triplet loss is added to the head-specific objective. The configured margin is zero, corresponding to the soft-margin formulation used in the FastReID implementation:

\begin{equation}
\mathcal{L}_{\mathrm{tri}}
=
\log\left(1+\exp(d_{ap}-d_{an})\right),
\end{equation}

where $d_{ap}$ and $d_{an}$ denote the Euclidean distances from an anchor to its hardest positive and hardest negative, respectively. These distances are computed using the unnormalized BNNeck features. Mining is performed within each device's sub-batch of 64 images, containing eight identities with eight images per identity.

The primary loss sweep comprises eight candidate configurations: linear softmax, ArcFace, Sub-center ArcFace, and SphereFace2, each evaluated both without and with hard-mined triplet loss. Thus, triplet supervision is an explicit comparison factor for every candidate head. The FastReID Circle-softmax-plus-triplet default is not part of this sweep; it is reserved as the reference objective in the separate comparison across the SBS, OSNet, and ConvNeXt-S. The eight candidates share a common \texttt{LossStudyBaseline} meta-architecture except where noted in Table~\ref{tab:conditions}.

\begin{table}[t]
\caption{Eight candidate objective configurations evaluated using the shared SBS R50-IBN architecture. Each head is trained with and without triplet loss.}
\label{tab:conditions}
\centering
\footnotesize
\begin{tabular}{@{}lll@{}}
\toprule
Classifier / head & Losses & Reference \\
\midrule
Linear softmax & CE & \cite{luo2019bag} \\
ArcFace ($s=64,m=0.35$) & CE & \cite{deng2019arcface} \\
Sub-center ArcFace ($K=3$) & CE & \cite{deng2020subcenter} \\
Angular head & SphereFace2 & \cite{wen2022sphereface2} \\
\midrule
Linear softmax & CE + Triplet & \cite{luo2019bag,hermans2017triplet} \\
ArcFace & CE + Triplet & \cite{deng2019arcface,hermans2017triplet} \\
Sub-center ArcFace ($K=3$) & CE + Triplet & \cite{deng2020subcenter,hermans2017triplet} \\
Angular head & SphereFace2 + Triplet & \cite{wen2022sphereface2,hermans2017triplet} \\
\bottomrule
\end{tabular}
\end{table}

\subsection{Evaluation Metrics}
\label{sec:metrics}

We evaluate retrieval performance using mean average precision (mAP) and the cumulative matching characteristic (CMC), following the standard query and gallery protocols of each benchmark. CMC performance is summarized using Rank-1 and Rank-5 accuracy, representing the proportion of queries for which a correct match appears within the top one and top five retrieved gallery images, respectively. Full CMC curves up to Rank-20 are additionally reported for selected models. Extracted features are $\ell_2$-normalized and ranked using cosine similarity. No re-ranking is applied. All models are evaluated using identical query and gallery roles within each benchmark, ensuring consistent comparisons across objective configurations and architectures.


%==============================================================================

\section{Results}
\label{sec:results}

The primary objective comparison comprises 24 training runs (eight candidate configurations $\times$ three seeds). Held-out retrieval is reported on four public benchmarks---Market-1501, MSMT17, DukeMTMC-reID, and Occluded-Duke---using the unseen-target evaluation bundle. Each table cell gives mean mAP / Rank-1 (\%) with seed standard deviations ($\pm$). The average column is the unweighted mean across the four benchmarks. Rank-5 accuracy and CMC curves provide additional analysis; the default-versus-selected cross-backbone comparison is reported separately.

\subsection{Objective Comparison}
\label{sec:objective-results}

Tables~\ref{tab:noTriplet} and~\ref{tab:triplet} compare the evaluated classification objectives without and with triplet supervision, respectively. This comparison examines their retrieval performance under the same SBS architecture and multi-source training protocol.

\begin{table}[t]
\caption{Objective comparison without triplet loss: mean mAP / Rank-1 (\%) over seeds 0--2 on four held-out benchmarks.}
\label{tab:noTriplet}
\centering
\scriptsize
\setlength{\tabcolsep}{3pt}
\begin{tabular}{@{}l*{5}{c}@{}}
\toprule
Objective & Market & MSMT17 & Duke & Occ-Duke & Avg. \\
\midrule
Linear softmax & \mstd{58.7}{0.5} / \mstd{81.0}{0.6} & \mstd{16.2}{0.3} / \mstd{36.3}{0.5} & \mstd{48.8}{0.4} / \mstd{65.7}{0.7} & \mstd{21.0}{0.7} / \mstd{26.9}{0.6} & \mstd{36.2}{0.2} / \mstd{52.5}{0.2} \\
ArcFace & \mstd{69.8}{0.3} / \mstd{87.8}{0.2} & \mstd{23.5}{0.4} / \mstd{49.3}{0.5} & \mstd{56.9}{0.5} / \mstd{73.4}{0.6} & \mstd{28.9}{0.5} / \mstd{37.1}{0.6} & \mstd{44.8}{0.3} / \mstd{61.9}{0.1} \\
Sub-center ArcFace & \TBD & \TBD & \TBD & \TBD & \TBD \\
SphereFace2 & \mstd{60.3}{1.4} / \mstd{83.5}{0.4} & \mstd{16.7}{1.0} / \mstd{37.9}{1.2} & \mstd{43.6}{4.3} / \mstd{65.6}{2.5} & \mstd{21.0}{2.3} / \mstd{31.9}{2.2} & \mstd{35.4}{2.2} / \mstd{54.7}{1.4} \\
\bottomrule
\end{tabular}
\end{table}

\begin{table}[t]
\caption{Objective comparison with triplet loss: mean mAP / Rank-1 (\%) over seeds 0--2 on four held-out benchmarks.}
\label{tab:triplet}
\centering
\scriptsize
\setlength{\tabcolsep}{3pt}
\begin{tabular}{@{}l*{5}{c}@{}}
\toprule
Objective & Market & MSMT17 & Duke & Occ-Duke & Avg. \\
\midrule
Linear softmax & \TBD & \TBD & \TBD & \TBD & \TBD \\
ArcFace & \mstd{70.4}{0.8} / \mstd{87.9}{0.4} & \mstd{22.0}{0.2} / \mstd{46.6}{0.4} & \mstd{56.3}{0.5} / \mstd{72.4}{0.6} & \mstd{28.7}{0.5} / \mstd{35.8}{0.5} & \mstd{44.3}{0.2} / \mstd{60.7}{0.1} \\
Sub-center ArcFace & \mstd{67.6}{0.6} / \mstd{87.1}{0.4} & \mstd{20.7}{0.4} / \mstd{44.4}{0.8} & \mstd{54.1}{0.5} / \mstd{70.3}{0.6} & \mstd{26.2}{0.8} / \mstd{33.3}{1.1} & \mstd{42.1}{0.5} / \mstd{58.8}{0.7} \\
SphereFace2  & \mstd{62.2}{2.0} / \mstd{84.4}{0.9} & \mstd{16.3}{0.8} / \mstd{37.8}{1.9} & \mstd{45.2}{0.7} / \mstd{66.4}{0.6} & \mstd{23.0}{0.8} / \mstd{34.9}{0.4} & \mstd{36.7}{0.7} / \mstd{55.9}{0.7} \\
\bottomrule
\end{tabular}
\end{table}

Tables~\ref{tab:noTriplet} and~\ref{tab:triplet} report held-out results for the eight candidate objectives, separated by whether triplet supervision is used. The Circle-softmax-plus-triplet FastReID default is excluded from these tables and appears only in the cross-backbone comparison. The best candidate will be selected using the composite internal validation protocol and compared with the default across all three backbones.

\subsection{Cross-Backbone Comparison}
\label{sec:backbone-results}

To examine whether the validation-selected candidate transfers beyond the primary architecture, we compare it with the FastReID Circle-softmax-plus-triplet default using SBS R50-IBN, OSNet, and ConvNeXt-S. Table~\ref{tab:backbone} summarizes retrieval on the same four held-out image benchmarks.

\begin{table}[t]
\caption{Cross-backbone comparison (mAP / Rank-1, \%; mean$\pm$std over seeds 0--2 on four held-out benchmarks).}
\label{tab:backbone}
\centering
\scriptsize
\setlength{\tabcolsep}{3pt}
\begin{tabular}{@{}ll*{5}{c}@{}}
\toprule
Backbone & Objective & Market & MSMT17 & Duke & Occ-Duke & Avg. \\
\midrule
SBS R50-IBN & Circle-softmax + Triplet & \mstd{67.4}{1.3} / \mstd{86.7}{0.6} & \mstd{15.5}{0.2} / \mstd{37.5}{0.4} & \mstd{52.5}{1.1} / \mstd{71.2}{1.1} & \mstd{26.9}{1.0} / \mstd{36.2}{1.2} & \mstd{40.6}{0.8} / \mstd{57.9}{0.6} \\
SBS R50-IBN & Selected candidate & \TBD & \TBD & \TBD & \TBD & \TBD \\
\midrule
OSNet & Circle-softmax + Triplet & \TBD & \TBD & \TBD & \TBD & \TBD \\
OSNet & Selected candidate & \TBD & \TBD & \TBD & \TBD & \TBD \\
\midrule
ConvNeXt-S & Circle-softmax + Triplet & \TBD & \TBD & \TBD & \TBD & \TBD \\
ConvNeXt-S & Selected candidate & \TBD & \TBD & \TBD & \TBD & \TBD \\
\bottomrule
\end{tabular}
\end{table}

This comparison tests whether the candidate selected by internal validation retains its relative performance against the FastReID default across architectures. \todo{Discuss the cross-backbone results and any architecture-dependent differences.}

\subsection{Retrieval Performance Analysis}
\label{sec:detailed}

In addition to mAP and Rank-1, Table~\ref{tab:rank5} reports Rank-5 accuracy for the eight candidate objective configurations. Figure~\ref{fig:cmc} compares the validation-selected candidate with the Circle-softmax-plus-triplet FastReID default across the public benchmarks, providing a view of retrieval performance beyond the first few ranks.

\begin{table}[t]
\caption{Rank-5 accuracy (\%): mean$\pm$std over seeds 0--2 for the eight candidate objective configurations on four held-out benchmarks.}
\label{tab:rank5}
\centering
\scriptsize
\setlength{\tabcolsep}{3pt}
\begin{tabular}{@{}l*{5}{c}@{}}
\toprule
Objective & Market & MSMT17 & Duke & Occ-Duke & Avg. \\
\midrule
Linear softmax (CE) & \mstd{89.9}{0.4} & \mstd{50.4}{0.5} & \mstd{79.1}{0.2} & \mstd{41.1}{1.0} & \mstd{65.1}{0.0} \\
ArcFace (CE) & \mstd{94.7}{0.1} & \mstd{63.0}{0.6} & \mstd{85.0}{0.3} & \mstd{51.4}{0.5} & \mstd{73.5}{0.1} \\
Sub-center ArcFace (CE) & \TBD & \TBD & \TBD & \TBD & \TBD \\
SphereFace2 & \mstd{93.5}{0.3} & \mstd{53.2}{1.0} & \mstd{79.1}{2.1} & \mstd{47.3}{2.8} & \mstd{68.3}{1.5} \\
Linear softmax + Triplet & \TBD & \TBD & \TBD & \TBD & \TBD \\
ArcFace + Triplet & \mstd{95.0}{0.4} & \mstd{61.0}{0.2} & \mstd{84.6}{0.4} & \mstd{51.6}{0.9} & \mstd{73.0}{0.3} \\
Sub-center ArcFace + Triplet & \mstd{94.4}{0.1} & \mstd{58.2}{1.0} & \mstd{83.2}{0.8} & \mstd{48.8}{1.7} & \mstd{71.2}{0.9} \\
SphereFace2 + Triplet & \mstd{93.5}{0.4} & \mstd{53.2}{1.7} & \mstd{80.1}{0.8} & \mstd{50.2}{0.9} & \mstd{69.3}{0.5} \\
\bottomrule
\end{tabular}
\end{table}

\begin{figure}[t]
\centering
\fbox{\parbox{0.92\textwidth}{\centering\vspace{1.5em}
\textbf{Figure 2 (placeholder).}\\
CMC curves (Ranks 1--20) across the public evaluation benchmarks for the validation-selected candidate and the FastReID Circle-softmax-plus-triplet default using SBS R50-IBN.
\vspace{1.5em}}}
\caption{CMC curves comparing the validation-selected candidate with the FastReID default on the public evaluation benchmarks.}
\label{fig:cmc}
\end{figure}

\todo{Discuss whether Rank-5 and the CMC curves support the trends observed in mAP and Rank-1, including any differences between benchmarks.}

\subsection{Computational Cost}
\label{sec:efficiency}

Table~\ref{tab:cost} compares the parameter counts, computational requirements, and inference speeds of the three convolutional architectures. These measurements provide additional context for interpreting their retrieval performance under practical computational constraints. For reference, published model-size figures for representative Transformer-based Re-ID approaches may also be included, provided that their measurement conditions and sources are clearly identified.

\begin{table}[t]
\caption{Model size and inference cost of the evaluated convolutional architectures.}
\label{tab:cost}
\centering
\scriptsize
\setlength{\tabcolsep}{3pt}
\begin{tabular}{@{}lccccc@{}}
\toprule
Model & Params (M) & FLOPs (G) & Latency (ms/img) & Throughput (img/s) \\
\midrule
SBS R50-IBN & \TBD & \TBD & \TBD & \TBD \\
OSNet & \TBD & \TBD & \TBD & \TBD \\
ConvNeXt-S & \TBD & \TBD & \TBD & \TBD \\
\bottomrule
\end{tabular}
\end{table}

Inference latency and throughput are measured under a consistent hardware and software configuration, using the same input resolution, batch size, and numerical precision. \todo{Specify the GPU, software environment, batch size, precision, and input resolution used for timing.}

%==============================================================================
\section{Discussion}
\label{sec:discussion}

\todo{Interpret Tables~\ref{tab:noTriplet}--\ref{tab:cost} without repeating every number.} The analysis will compare the objective families and the effect of triplet loss, then examine whether performance differences vary across benchmarks with different scales and occlusion levels and carry over to Alexandria. It will also consider whether observed trends are qualitatively consistent with the heads' assumptions about noisy or multi-modal classes; because no controlled label-noise ablation is performed, this interpretation will not be treated as causal. The cross-backbone comparison tests whether the validation-selected candidate transfers to OSNet and ConvNeXt-S relative to the Circle-softmax-plus-triplet default. All reported differences will be interpreted alongside variation across the three seeds, without formal significance claims.

%==============================================================================
\section{Conclusion}
\label{sec:conclusion}

This work examines how identity-learning objectives affect cross-domain generalization in person Re-ID. Eight candidate configurations, covering four identity-learning heads with and without triplet supervision, are compared on the FastReID Strong Baseline. The strongest candidate selected by composite internal validation is then compared with the Circle-softmax-plus-triplet FastReID default across the Strong Baseline, OSNet, and ConvNeXt-S. \todo{Summarize the selected candidate's gains and practical implications after results are finalized.} Completion of the cross-source leakage review, extension and public release of the Alexandria dataset, and evaluation on transformer-based backbones remain directions for future work.

%==============================================================================
\section{Limitations}
\label{sec:limitations}

The primary objective comparison covers eight candidate configurations on the ImageNet-pretrained ResNet-50-IBN Strong Baseline under a fixed fifty-epoch budget; its conclusions may not extend to other heads or architectures. The Circle-softmax-plus-triplet FastReID default is evaluated separately against the candidate selected using internal validation, with the cross-backbone comparison on OSNet and ConvNeXt-S underway and numerical values pending. Internal validation is a pooled identity-holdout proxy rather than a public benchmark, so selected checkpoints may not coincide with the optimum on held-out benchmarks. In addition, the cross-source pixel-level duplicate check across the nine training sources and test sets is not yet complete, so strict acquisition independence cannot be guaranteed; DukeMTMC and Occluded-Duke share acquisition history. The study contains no controlled label-noise or identity-multimodality ablation, so interpretations of those properties are qualitative.

Finally, the availability of DukeMTMC depends on licensing, and release of the Alexandria data depends on institutional approval.

\begin{credits}
% \subsubsection{\ackname}
% We thank the Faculty of Computers and Data Science, Alexandria University, for supporting data collection and compute.

\subsubsection{\discintname}
The authors have no competing interests to declare that are relevant to the content of this article.
\end{credits}

\bibliographystyle{splncs04}
\bibliography{references}

\end{document}